\begin{explanationbox}
\textbf{\textcolor{CExplanation}{一句话直觉}}
直线与椭圆相切等价于联立方程后二次判别式为零，导出简洁的 $m^2$ 表达式。

\medskip
\textbf{\textcolor{CExplanation}{核心拆解}}
\begin{description}[style=nextline, leftmargin=2.8em, labelsep=0.5em, itemsep=0.45em]
    \item[\textbf{要点一（联立方程）：}] 将直线代入椭圆，得到关于 $x$ 的二次方程。
    \item[\textbf{要点二（判别式零）：}] 相切时二次方程有重根，故判别式 $\Delta = 0$。
\end{description}

\medskip
\textbf{\textcolor{CExplanation}{几何本质（唯一交点）}}
切线与椭圆仅有一个公共点，且该点处直线斜率与椭圆局部斜率一致，但代数上只需用公共点唯一性刻画。

\medskip
\textbf{\textcolor{CExplanation}{代数意义（重根条件）}}
联立后的二次方程 $Ax^2+Bx+C=0$ 有相等实根，因此判别式 $B^2-4AC=0$，化简即得 $m^2=a^2k^2+b^2$。

\medskip
\textbf{\textcolor{CExplanation}{考点价值（切线问题）}}
\begin{itemize}[leftmargin=*, itemsep=0.5em, topsep=0.4em]
    \item \textbf{考法一（求切线）：} 已知斜率 $k$，直接套公式写出切线方程。
    \item \textbf{考法二（求参数）：} 已知直线与椭圆相切，反求 $k$ 或 $m$ 的值。
\end{itemize}

\medskip
\textbf{\textcolor{CExplanation}{顿悟点}}
\textbf{$m^2$ 的表达式中 $a^2k^2$ 与 $b^2$ 相加，形式对称，便于记忆；它本质是判别式为零的紧凑形式。}

\medskip
\textbf{\textcolor{CExplanation}{使用场景}}
\begin{itemize}[leftmargin=*, itemsep=0.5em, topsep=0.4em]
    \item \textbf{场景一（判断相切）：} 给定直线与椭圆，快速验证是否相切。
    \item \textbf{场景二（写切线）：} 已知椭圆和斜率，直接写出切线方程。
\end{itemize}
\end{explanationbox}
